% fig_vl11_kasten.tex — Teilchen im Kasten: Eigenfunktionen auf ihren
% Energieniveaus, und die quadratische Aufspreizung des Spektrums.
%
% Reference figure for the notes: the other lecture figures follow the
% conventions established here.
%
%   * geometry and curves come from the closed-form solution derived in
%     the lecture, never from a freehand sketch:
%         psi_n(x) = sqrt(2/L) sin(n pi x / L),
%         E_n      = n^2 pi^2 hbar^2 / (2 m L^2),
%     so the node count, the symmetry and the level spacing are forced
%     to be right rather than drawn by eye;
%   * each eigenfunction is drawn about its own energy — the standard way
%     of showing that psi_n belongs to E_n;
%   * colour encodes the quantum number n (the ordered ramp of
%     `tpclevels`), and nothing else.  The energy levels themselves are
%     neutral grey rules: they are reference lines, not data;
%   * the classically forbidden region is hatched, not filled, so it
%     survives greyscale printing.
%
% Panel (a) the first four eigenfunctions in the well.
% Panel (b) the spectrum E_n/E_1, with the oscillator's evenly spaced
%           ladder alongside it — the contrast between a quadratic and a
%           linear spectrum is the point students should carry away.

\begin{center}
\begin{tikzpicture}

% ─────────────────────────────────────────────────────────────
%  (a)  Eigenfunctions drawn about their levels
% ─────────────────────────────────────────────────────────────
\begin{axis}[
    name=wells,
    tpcplot,
    width=8.2cm, height=6.8cm,
    axis lines=none,
    xmin=-0.34, xmax=1.34,
    ymin=-1.2,  ymax=19.4,
    clip=false,
    title={(a) \tpcifenglish{$\psi_n$ on their levels}%
                       {$\psi_n$ auf ihren Niveaus}},
    domain=0:1, samples=201,
]

  % Classically forbidden regions, x<0 and x>L, where V = infinity
  \addplot[forbidden, draw=none, domain=-0.34:0, samples=2] {19.4} \closedcycle;
  \addplot[forbidden, draw=none, domain=1:1.34,  samples=2] {19.4} \closedcycle;

  % The well itself
  \draw[potential] (axis cs:0,-1.2) -- (axis cs:0,19.4);
  \draw[potential] (axis cs:1,-1.2) -- (axis cs:1,19.4);
  \draw[potential] (axis cs:0,0)    -- (axis cs:1,0);

  % Energy levels E_n = n^2 E_1: neutral reference rules
  \addplot[levellinegrey, domain=0:1, samples=2] {1};
  \addplot[levellinegrey, domain=0:1, samples=2] {4};
  \addplot[levellinegrey, domain=0:1, samples=2] {9};
  \addplot[levellinegrey, domain=0:1, samples=2] {16};

  % Eigenfunctions about those levels.  Amplitude 1.3 is the largest that
  % keeps psi_2 clear of E_1, the tightest gap (E_2 - E_1 = 3 E_1).
  \addplot[wavefn, spacecadet]   {1  + 1.3*sin(deg(1*pi*x))};
  \addplot[wavefn, cgblue]       {4  + 1.3*sin(deg(2*pi*x))};
  \addplot[wavefn, munsell]      {9  + 1.3*sin(deg(3*pi*x))};
  \addplot[wavefn, birrenorange] {16 + 1.3*sin(deg(4*pi*x))};

  % Level labels.  Written out rather than looped: \foreach defers
  % expansion inside a pgfplots axis.
  \node[lbl, anchor=east] at (axis cs:-0.06,1)  {$E_{1}$};
  \node[lbl, anchor=east] at (axis cs:-0.06,4)  {$E_{2}$};
  \node[lbl, anchor=east] at (axis cs:-0.06,9)  {$E_{3}$};
  \node[lbl, anchor=east] at (axis cs:-0.06,16) {$E_{4}$};

  % Node counts: psi_n has n-1 interior nodes
  \node[lbl, text=birrengray, font=\scriptsize, anchor=west]
    at (axis cs:1.03,1)  {$n=1$};
  \node[lbl, text=birrengray, font=\scriptsize, anchor=west]
    at (axis cs:1.03,4)  {$n=2$};
  \node[lbl, text=birrengray, font=\scriptsize, anchor=west]
    at (axis cs:1.03,9)  {$n=3$};
  \node[lbl, text=birrengray, font=\scriptsize, anchor=west]
    at (axis cs:1.03,16) {$n=4$};

  % Axis furniture
  \node[lbl, anchor=north] at (axis cs:0,-0.3)    {$0$};
  \node[lbl, anchor=north] at (axis cs:1,-0.3)    {$L$};
  \node[lbl, anchor=north] at (axis cs:0.5,-0.62) {$x$};
  \node[lbl, text=birrengray, font=\scriptsize, anchor=south]
    at (axis cs:-0.17,17.6) {$V=\infty$};
\end{axis}

% ─────────────────────────────────────────────────────────────
%  (b)  Spectrum: the n^2 ladder against the oscillator's linear one
% ─────────────────────────────────────────────────────────────
\begin{axis}[
    tpcplot,
    at={(wells.south east)}, anchor=south west, xshift=17mm,
    width=6.2cm, height=6.8cm,
    xlabel={$n$},
    ylabel={$E_n/E_1$},
    every axis y label/.style={at={(axis description cs:-0.16,0.5)},
                               rotate=90, anchor=south, font=\small},
    xmin=0.4, xmax=5.6, ymin=0, ymax=26.5,
    xtick={1,2,3,4,5}, ytick={0,5,10,15,20,25},
    grid=major,
    legend pos=north west,
    legend style={font=\scriptsize, draw=light-gray},
    title={(b) \tpcifenglish{Spectrum}{Spektrum}},
]
  % Kasten: E_n / E_1 = n^2
  \addplot[cgblue, line width=1.0pt, mark=*, mark size=1.9pt,
           domain=1:5, samples=5] {x^2};
  \addlegendentry{\tpcifenglish{box}{Kasten}, $n^2$}

  % Oscillator, normalised the same way: E_n/E_1 = (n + 1/2)/(3/2)
  \addplot[birrenorange, line width=1.0pt, densely dashed,
           mark=o, mark size=1.9pt, mark options={solid},
           domain=1:5, samples=5] {(2*x+1)/3};
  \addlegendentry{\tpcifenglish{oscillator}{Oszillator}, $\tfrac{2n+1}{3}$}
\end{axis}

\end{tikzpicture}
\end{center}
